% Dataset & EDA section (Part B)
% Include from reports/midsem/main.tex.

\section{Dataset and Exploratory Analysis}
\label{sec:dataset-eda}

\subsection{EuRoC MAV}
We use three Machine Hall sequences from the EuRoC MAV dataset
\cite{burri2016euroc}: MH\_01\_easy, MH\_02\_easy, and MH\_03\_medium.
Each flight is a few minutes of a real hexacopter with a stereo camera and
an ADIS16448 IMU. Phase~1 uses only \texttt{imu0/data.csv} (gyroscope in
rad/s, accelerometer specific force in m/s$^2$, nominally 200\,Hz) and
\texttt{state\_groundtruth\_estimate0/data.csv} (position, Hamilton
quaternion, velocity, and IMU biases). Camera images are reserved for
Phase~2.

\subsection{Data quality}
Table~\ref{tab:quality} summarises the three streams after loading.
There are no missing values. Both IMU and the provided ground-truth
estimate are sampled at 200\,Hz with a 5\,ms median step and negligible
jitter. Nearest-timestamp alignment drops 1.17\%, 1.32\%, and 2.58\% of
IMU rows (MH\_01--MH\_03), almost entirely at sequence edges where
ground-truth coverage starts later than the IMU. We treat those drops as
coverage gaps, not sensor faults.

\begin{table}[t]
\centering
\caption{Sequence-level quality checks.}
\label{tab:quality}
\begin{tabular}{lrrrrr}
\hline
Sequence & IMU rows & Duration (s) & Rate (Hz) & Missing & Align drop \\
\hline
MH\_01\_easy & 36820 & 184.1 & 200 & 0 & 1.17\% \\
MH\_02\_easy & 30400 & 152.0 & 200 & 0 & 1.32\% \\
MH\_03\_medium & 27008 & 135.0 & 200 & 0 & 2.58\% \\
\hline
\end{tabular}
\end{table}

Raw gyro over the first 10\,s of MH\_01 is small-amplitude chatter
($\sim$0.2\,rad/s), i.e.\ noise plus slow attitude change rather than
aggressive spins. Accelerometer magnitude stays near $g$; the vehicle is
pitched, so gravity appears on $a_x$ and $a_z$ together, not on $a_z$
alone.

\subsection{Easy versus medium}
Ground-truth speed distinguishes the labels: mean 0.44\,m/s and
0.49\,m/s on the two easy flights versus 0.99\,m/s on MH\_03 (peaks
2.8\,m/s). Path length is $\sim$81\,m and $\sim$74\,m versus $\sim$131\,m.
Mean gyro magnitude is also higher on MH\_03 (0.29 vs 0.21\,rad/s).
The 3-D trajectories occupy a similar machine-hall volume; MH\_03 is
``medium'' because of speed and turning, not a larger room.

\subsection{Drift label}
The headline target is the 1\,s raw-IMU position error $y=\texttt{err\_mag\_m}$.
Across 724 / 596 / 523 overlapping windows the means are 0.186, 0.186,
and 0.181\,m with standard deviations 0.012, 0.013, and 0.021\,m.
The easy-sequence histograms are tight and slightly left-skewed; MH\_03
is a little wider (max 0.30\,m) but still not a long tail. We therefore
\textbf{do not} log-transform the primary target.

The same windows after subtracting the dataset's ground-truth IMU biases
(\texttt{err\_mag\_m\_bc}) have means 0.021, 0.018, and 0.030\,m. MH\_03
then shows the most residual structure (max 0.176\,m, skew 2.3). This
supports the claim that 1\,s raw drift is bias-dominated; the
bias-corrected series is an oracle comparison, not the leaderboard
target. Signed axis errors $({\Delta}x,{\Delta}y,{\Delta}z)$ vary more
than the Euclidean magnitude, so Phase~1 models should also be trained
on the 3-axis target.

Mean error versus window length (0.25--2.0\,s) grows superlinearly for
the raw integrator ($\sim$6\,mm at 0.25\,s to $\sim$1.2\,m at 2\,s) and
remains much smaller after bias correction. The primary window stays
1.0\,s / 0.25\,s stride as specified.

\subsection{Train/test protocol}
Windows overlap by 75\%. A random split of rows would place nearly
identical IMU slices in both train and test and inflate $R^2$. We use
leave-one-sequence-out, recorded in \texttt{data/processed/splits.json}:
train on two flights, test on the third, rotated three times. The
leakage test \texttt{tests/test\_no\_leakage.py} asserts that no
sequence appears on both sides of any fold. The primary target in that
file is \texttt{err\_mag\_m}.
